\newpage

\renewcommand\thesection{3}
\renewcommand\thesubsection{\arabic{subsection}}

\counterwithin{subsection}{section}
\setcounter{subsection}{0}

\counterwithin{table}{section}
\captionsetup[table]{labelsep=quad}
\renewcommand{\thetable}{\arabic{section}.\arabic{table}}

\begin{center}
    \section*{\textbf{BAB III \\ METODE PELAKSANAAN}}
    \addcontentsline{toc}{section}{BAB III \quad METODE PELAKSANAAN}
\end{center}

\subsection{Metode penelitian}

\hspace{0.5cm}Metode penelitian menggunakan \textit{mixed method research} yang mengombinasikan 2 metode yakni metode kualitatif (\textit{qualitative method}) dan metode kuantitatif (\textit{quantitative method}) \citep{subedi2023}. Metode kualitatif yang digunakan dalam penelitian ini berupa wawancara semi-terstruktur dengan lima tokoh agama dari latar belakang agama yang berbeda, yaitu agama Buddha, Kristen, Katolik, Hindu, dan Islam. Metode kuantitatif penelitian ini menggunakan data survei kuesioner yang disebarkan kepada pelajar, mahasiswa, dan masyarakat umum untuk mengetahui pandangan responden mengenai hubungan antara nilai agama, moralitas, dan sikap antikorupsi.

\subsection{Lokasi penelitian}

Wawancara narasumber dilakukan di beberapa lokasi kegiatan yang berbeda sesuai dengan latar belakang agama narasumber. Wawancara agama Buddha dilakukan bersama Aryabodhi Heru Budi Santoso di Vihara Ekayana, wawancara agama Kristen dilakukan bersama Kevin Manuelim di GPKdI Blessing Community Jelambar, wawancara agama Katolik dilakukan bersama Romo Greg Soetomo SJ di Komunitas Mahasiswa Katolik Jakarta Barat, wawancara agama Hindu dilakukan bersama pihak Pasraman Chandra Prabha, dan wawancara agama Islam dilakukan bersama Ustadz Masruchin SHI sebagai Penyuluh Agama Islam Kec. Grogol Petamburan. Pemilihan narasumber dari lima agama dilakukan untuk memperoleh pandangan yang lebih luas mengenai bagaimana nilai keagamaan dapat berperan dalam menanamkan moral antikorupsi kepada masyarakat.

\subsubsection{Waktu penelitian}
Penelitian dilakukan mulai dari tahap perencanaan, penyusunan instrumen, pelaksanaan wawancara, penyebaran kuesioner, hingga analisis data. Kunjungan wawancara dilakukan pada tanggal 15 Maret 2026 hingga 11 April 2026 secara onsite di lokasi kegiatan masing-masing narasumber. Pengambilan data survei menggunakan kuesioner dilakukan pada periode 15 April hingga 11 Mei 2026 dan disebarkan secara langsung maupun daring melalui media sosial dan platform komunikasi digital.

\subsubsection{Teknik pengumpulan data dan sampling}

Pengumpulan data dalam penelitian ini dilakukan melalui dua teknik, yaitu wawancara dan kuesioner. Wawancara dilakukan secara semi-terstruktur kepada lima tokoh agama yang dipilih berdasarkan keterkaitan narasumber dengan topik nilai keagamaan dan pembinaan moral masyarakat. Sementara itu, kuesioner disebarkan kepada pelajar, mahasiswa, dan masyarakat umum untuk memperoleh data kuantitatif mengenai pandangan responden terhadap agama, korupsi, serta peran pendidikan moral antikorupsi. Teknik sampling untuk wawancara menggunakan \textit{purposive sampling}, sedangkan pengambilan responden kuesioner dilakukan melalui penyebaran terbuka kepada responden yang bersedia mengisi kuesioner.

\subsection{Analisis data}

Data yang didapatkan dari kuesioner dikumpulkan, dibersihkan (\textit{cleaning}), dan dipersiapkan untuk dianalisis secara deskriptif. Analisis kuantitatif dilakukan melalui perhitungan frekuensi, persentase, dan rata-rata untuk menggambarkan kecenderungan jawaban responden terhadap setiap pertanyaan. Kuesioner terdiri dari pertanyaan ya atau tidak serta pertanyaan dengan skala Likert lima poin. Skala Likert digunakan untuk mengukur tingkat persetujuan responden terhadap pernyataan mengenai pandangan agama terhadap korupsi, peran tokoh agama, pendidikan moral, serta sikap antikorupsi dalam kehidupan sehari-hari \citep{alkharusi2022}. Skala Likert yang digunakan dapat dilihat pada Tabel \ref{tab:TabelSkalalikert}.
\stepcounter{section}
\stepcounter{section}
\stepcounter{section}
\begin{table}[h]
    \centering
    \vspace{10pt} % Tambahkan spasi sebelum judul tabel
    \captionsetup{labelformat=empty}
    \caption[Skala Likert]{Tabel \ref{tab:TabelSkalalikert} Skala Likert}
    \begin{tabular}{ |c|c| }
        \hline
        \textbf{Opsi pertanyaan} & \textbf{Penilaian} \\
        \hline
        Sangat setuju & 5 \\
        \hline
        Setuju & 4 \\
        \hline
        Netral & 3 \\
        \hline
        Tidak setuju & 2 \\
        \hline
        Sangat tidak setuju & 1 \\
        \hline
    \end{tabular}
    \label{tab:TabelSkalalikert}
\end{table}

Di dalam kuesioner terdapat beberapa pertanyaan yang dibagi berdasarkan aspek yang ingin diukur. Pertanyaan awal digunakan untuk mengetahui pengetahuan dasar responden mengenai larangan korupsi dalam agama, pengalaman mendengar ceramah antikorupsi, dan pengetahuan mengenai sanksi moral atau spiritual bagi pelaku korupsi. Selanjutnya, pertanyaan skala Likert digunakan untuk mengukur tingkat persetujuan responden terhadap pernyataan mengenai nilai agama, peran tokoh agama, pengaruh ceramah antikorupsi, pendidikan agama sejak kecil, serta konsekuensi moral dan spiritual dari tindakan korupsi.

Data kualitatif dari hasil wawancara dianalisis menggunakan teknik \textit{thematic analysis}. Analisis ini dilakukan dengan mengidentifikasi, mengelompokkan, dan menafsirkan tema-tema utama yang muncul dari jawaban narasumber \citep{christou2023}. Tema utama yang dianalisis meliputi pandangan masing-masing agama terhadap korupsi, faktor penyebab korupsi tetap terjadi di tengah masyarakat religius, serta peran lembaga keagamaan dan tokoh agama dalam membangun moral antikorupsi. Hasil analisis kuantitatif dan kualitatif kemudian dikombinasikan untuk memperoleh pemahaman yang lebih menyeluruh mengenai peran nilai keagamaan dalam penanaman moral antikorupsi kepada masyarakat.
